%% bare_conf.tex
%% V1.4b
%% 2015/08/26
%% by Michael Shell
%% See:
%% http://www.michaelshell.org/
%% for current contact information.
%%
%% This is a skeleton file demonstrating the use of IEEEtran.cls
%% (requires IEEEtran.cls version 1.8b or later) with an IEEE
%% conference paper.
%%
%% Support sites:
%% http://www.michaelshell.org/tex/ieeetran/
%% http://www.ctan.org/pkg/ieeetran
%% and
%% http://www.ieee.org/

%%*************************************************************************
%% Legal Notice:
%% This code is offered as-is without any warranty either expressed or
%% implied; without even the implied warranty of MERCHANTABILITY or
%% FITNESS FOR A PARTICULAR PURPOSE!
%% User assumes all risk.
%% In no event shall the IEEE or any contributor to this code be liable for
%% any damages or losses, including, but not limited to, incidental,
%% consequential, or any other damages, resulting from the use or misuse
%% of any information contained here.
%%
%% All comments are the opinions of their respective authors and are not
%% necessarily endorsed by the IEEE.
%%
%% This work is distributed under the LaTeX Project Public License (LPPL)
%% ( http://www.latex-project.org/ ) version 1.3, and may be freely used,
%% distributed and modified. A copy of the LPPL, version 1.3, is included
%% in the base LaTeX documentation of all distributions of LaTeX released
%% 2003/12/01 or later.
%% Retain all contribution notices and credits.
%% ** Modified files should be clearly indicated as such, including  **
%% ** renaming them and changing author support contact information. **
%%*************************************************************************


% *** Authors should verify (and, if needed, correct) their LaTeX system  ***
% *** with the testflow diagnostic prior to trusting their LaTeX platform ***
% *** with production work. The IEEE's font choices and paper sizes can   ***
% *** trigger bugs that do not appear when using other class files.       ***
% The testflow support page is at:
% http://www.michaelshell.org/tex/testflow/

\documentclass[conference]{IEEEtran}
% Some Computer Society conferences also require the compsoc mode option,
% but others use the standard conference format.
%
% If IEEEtran.cls has not been installed into the LaTeX system files,
% manually specify the path to it like:
% \documentclass[conference]{../sty/IEEEtran}

\usepackage{booktabs}
\usepackage{amsmath}
\usepackage{amsfonts}
\usepackage{amssymb}
\usepackage{graphicx}
\usepackage{float}
\usepackage{caption}
\usepackage{subcaption}
\usepackage[dvipsnames]{xcolor}
\usepackage{titlesec}

% ── Blue centred section titles ───────────────────────────────────────────────
\definecolor{ieeblue}{RGB}{0, 84, 166}

% Section: centred, large, bold, blue
\titleformat{\section}
  {\normalfont\large\bfseries\color{ieeblue}\centering}
  {\thesection.}{0.5em}{}

% Subsection: left-aligned, normalsize, bold, lighter blue
\titleformat{\subsection}
  {\normalfont\normalsize\bfseries\color{ieeblue!80}}
  {\thesubsection}{0.5em}{}
% ─────────────────────────────────────────────────────────────────────────────


% *** CITATION PACKAGES ***
%
%\usepackage{cite}
% cite.sty was written by Donald Arseneau
% V1.6 and later of IEEEtran pre-defines the format of the cite.sty package
% \cite{} output to follow that of the IEEE. Loading the cite package will
% result in citation numbers being automatically sorted and properly
% "compressed/ranged". e.g., [1], [9], [2], [7], [5], [6] without using
% cite.sty will become [1], [2], [5]--[7], [9] using cite.sty. cite.sty's
% \cite will automatically add leading space, if needed. Use cite.sty's
% noadjust option (cite.sty V3.8 and later) if you want to turn this off
% such as if a citation ever needs to be enclosed in parenthesis.
% cite.sty is already installed on most LaTeX systems. Be sure and use
% version 5.0 (2009-03-20) and later if using hyperref.sty.
% The latest version can be obtained at:
% http://www.ctan.org/pkg/cite
% The documentation is contained in the cite.sty file itself.


% *** GRAPHICS RELATED PACKAGES ***
%
\ifCLASSINFOpdf
  % \usepackage[pdftex]{graphicx}
  % declare the path(s) where your graphic files are
  % \graphicspath{{../pdf/}{../jpeg/}}
  % and their extensions so you won't have to specify these with
  % every instance of \includegraphics
  % \DeclareGraphicsExtensions{.pdf,.jpeg,.png}
\else
  % \usepackage[dvips]{graphicx}
  % \graphicspath{{../eps/}}
  % \DeclareGraphicsExtensions{.eps}
\fi


% correct bad hyphenation here
\hyphenation{op-tical net-works semi-conduc-tor}


\begin{document}

%
% paper title
% Titles are generally capitalized except for words such as a, an, and, as,
% at, but, by, for, in, nor, of, on, or, the, to and up, which are usually
% not capitalized unless they are the first or last word of the title.
% Linebreaks \\ can be used within to get better formatting as desired.
% Do not put math or special symbols in the title.
\title{Treatment-Aware Temporal Vision Transformer for Longitudinal Glioma Volume Prediction and Progression Video Generation}


% author names and affiliations
% use a multiple column layout for up to three different affiliations
\author{\IEEEauthorblockN{Mohamed Boufafa, Belkacem Chikhaoui, Belkacem Khaldi}
\IEEEauthorblockA{Applied Artificial Intelligence Institute\\
TELUQ University, Montreal, Canada \\
belkacem.chikhaoui@teluq.ca
}
}


% make the title area
\maketitle

% As a general rule, do not put math, special symbols or citations
% in the abstract
\begin{abstract}
Tracking how a brain tumor changes over time --- and why --- is one of the most
clinically urgent problems in neuro-oncology, yet most deep learning tools still
treat each MRI scan in isolation. We present a modular, three-stage pipeline
that reasons across a patient's entire longitudinal MRI history, incorporating
the therapies they received and the molecular profile of their tumor. Applied to
the MU-Glioma cohort (203~patients, 596~scans), our system comprises:
(1)~a fine-tuned nnU-Net backbone (30.8M~parameters) that segments whole tumor,
tumor core, and enhancing tumor with a mean Dice of~0.877;
(2)~TaViT, a Treatment-Aware temporal Vision Transformer (3.06M~parameters)
conditioned on per-scan treatment tokens and 26 molecular biomarkers, which
estimates tumor subregion volume trajectories across all observed timepoints,
reaching $R^2 = 0.908/0.953/0.952$ for the three subregions; and
(3)~a FiLM-conditioned DDPM (24.7M~parameters) that generates smooth,
anatomically consistent segmentation frames between consecutive scan pairs,
turning discrete measurements into a continuous tumor evolution video.
External zero-shot evaluation on BraTS-PTG (559~patients, 1,620~scans) confirms
generalization across sites. Ablation studies show that treatment conditioning
alone drives a~$+26.3\%$ improvement in trajectory accuracy.
\end{abstract}

\IEEEpeerreviewmaketitle

% ============================================================================
\section{Introduction}
% ============================================================================

Gliomas are the most common malignant primary brain tumors, accounting for
roughly 80\% of cases~\cite{stupp2005glioblastoma, osborn2022who}.
Glioblastoma, the most aggressive form, carries a median survival of under
15~months even with surgery, radiotherapy, and temozolomide
chemotherapy~\cite{stupp2005glioblastoma} --- and recurrence is nearly
inevitable. To track disease course, clinicians rely on serial MRI scans
assessed through the RANO criteria~\cite{wen2010rano}, yet in practice this
review remains largely manual and rarely incorporates the full treatment and
molecular context documented in the patient record.

Over the past decade, deep learning has transformed brain tumor segmentation.
Tools such as nnU-Net~\cite{isensee2021nnunet} and
SwinUNETR~\cite{tang2022swinunetr} can now delineate tumor subregions with
radiologist-level accuracy --- but only on a single scan at a time. They tell
us \emph{where} the tumor is, not \emph{how it has been changing}. Physics-based
reaction-diffusion models attempt to fill this gap but require manual
calibration and ignore the effect of treatment entirely~\cite{liu2023tadiff}.
GAN-based approaches can predict pre-operative tumor shapes~\cite{liu2023tadiff},
but again without any treatment conditioning. The most closely related work,
TaDiff~\cite{liu2023tadiff}, is a meaningful step forward: it conditions a
diffusion model on past MRI sequences and treatment history to synthesize
realistic future scan appearances. However, TaDiff was trained on just
18 patients and tested on 5, it does not report volumetric trajectory metrics,
and it does not incorporate patient molecular biomarkers such as IDH or MGMT
status.

In this paper we ask a complementary question: given a patient's full
longitudinal MRI history, their treatment record, and their molecular profile,
can we \emph{continuously estimate how tumor subregion volumes evolve} across
all observed timepoints --- and turn those estimates into an interpretable video
of tumor progression? Our three contributions are:

\begin{enumerate}

  \item \textbf{TaViT --- Treatment-Aware Volume Trajectory Transformer:}
    A bidirectional Vision Transformer (3.06M parameters) that sees a patient's
    entire scan sequence at once, conditioned on 8-dimensional per-scan
    treatment tokens and 26-dimensional molecular biomarkers (IDH1, MGMT, EGFR,
    and others), and estimates whole-tumor, tumor-core, and enhancing-tumor
    volumes at every timepoint simultaneously. It achieves
    $R^2 = 0.908/0.953/0.952$ across the three subregions.

  \item \textbf{A segmentation-pair DDPM for progression visualization:}
    A FiLM-conditioned~\cite{perez2018film} diffusion model~\cite{ho2020ddpm}
    that takes two consecutive segmentation maps as boundary conditions and
    generates smooth, anatomically consistent intermediate frames, producing a
    per-patient tumor evolution video with BraTS hierarchy enforced
    ($\text{ET} \subseteq \text{TC} \subseteq \text{WT}$).

  \item \textbf{Large-scale, multi-site validation:}
    We train on MU-Glioma (203~patients, 596~scans --- nearly $9\times$ the
    size of TaDiff's training set) and validate zero-shot on
    BraTS-PTG~\cite{deverdier2024bratsptg} (559~unseen patients, 1,620~scans
    with expert ground-truth segmentations), using an 18-test battery covering
    morphological, temporal, and clinical metrics.

\end{enumerate}


% ============================================================================
\section{Related Work}
% ============================================================================



% ============================================================================
\section{Methodology}
\label{sec:methodology}
% ============================================================================


\begin{figure*}[h]
\centering
% TODO: Insert pipeline overview diagram
\fbox{\parbox{0.95\textwidth}{\centering\vspace{2em}\textbf{[Pipeline Overview Diagram]}\\
\small MRI $\to$ nnU-Net Segmentation $\to$ 2,825-D Embeddings $\to$ TaViT (+ Treatment + Molecular) $\to$ Volume Predictions \& DDPM $\to$ Video\vspace{2em}}}
\caption{Overview of the proposed pipeline.}
\label{fig:pipeline}
\end{figure*}



% ============================================================================
\section{Experiments}
\label{sec:experiments}
% ============================================================================

\subsection{Datasets}


\subsection{Implementation}


\subsection{Results}


% ============================================================================
\section{Discussion}
\label{sec:discussion}
% ============================================================================



% ============================================================================
\section{Conclusion}
\label{sec:conclusion}
% ============================================================================



\bibliographystyle{IEEEtran}
\bibliography{ref}
\end{document}
